\documentclass{article}
\usepackage{amsmath,amssymb,amsthm}
\usepackage{algorithm,algorithmic}
\usepackage{bm}
\usepackage{hyperref}
\newtheorem{theorem}{Theorem}
\newtheorem{lemma}{Lemma}
\newtheorem{assumption}{Assumption}
\newtheorem{corollary}{Corollary}
\newcommand{\EE}{\mathbb{E}}
\newcommand{\RR}{\mathbb{R}}
\newcommand{\norm}[1]{\left\|#1\right\|}
\newcommand{\ip}[2]{\langle #1,#2\rangle}
\begin{document}

\section*{Algorithm Name}
\textbf{Stochastic Subsampled Newton (SSN)} for non‑convex smooth optimization.

\section*{Mathematical Formulation}
Let $f:\RR^{d}\rightarrow\RR$ be twice continuously differentiable.
At iteration $t$ we have a stochastic gradient $g_t$ and a stochastic
Hessian estimate $H_t\in\RR^{d\times d}$ obtained by uniform subsampling
of a mini‑batch of data points.  The update rule is
\begin{equation}
\boxed{
x_{t+1}=x_t-\alpha\bigl(H_t+\lambda I\bigr)^{-1}g_t},
\qquad t=0,1,\dots ,
\label{eq:update}
\end{equation}
where
\begin{itemize}
\item $\alpha>0$ is a (possibly diminishing) stepsize,
\item $\lambda>0$ is a regularisation parameter ensuring that $H_t+\lambda I$
  is positive definite and thus invertible.
\end{itemize}
The stochastic objects satisfy
\[
\EE[g_t\mid x_t]=\nabla f(x_t),\qquad
\EE[H_t\mid x_t]=\nabla^2 f(x_t).
\]

\section*{Pseudocode}
\begin{algorithm}[H]
\caption{Stochastic Subsampled Newton (SSN)}
\begin{algorithmic}[1]
\STATE \textbf{Input:} stepsize $\alpha>0$, regulariser $\lambda>0$, 
batch size $b$, max iterations $T$.
\STATE Initialise $x_0\in\RR^{d}$.
\FOR{$t=0$ \TO $T-1$}
    \STATE Sample a mini‑batch $\mathcal{B}_t$ of size $b$.
    \STATE Compute stochastic gradient
          $g_t=\frac{1}{b}\sum_{i\in\mathcal{B}_t}\nabla f_i(x_t)$.
    \STATE Compute stochastic Hessian
          $H_t=\frac{1}{b}\sum_{i\in\mathcal{B}_t}\nabla^2 f_i(x_t)$.
    \STATE $x_{t+1}=x_t-\alpha\bigl(H_t+\lambda I\bigr)^{-1}g_t$.
\ENDFOR
\STATE \textbf{Output:} $\{x_t\}_{t=0}^{T}$ (or the averaged iterate
$\bar x_T=\frac1T\sum_{t=0}^{T-1}x_t$).
\end{algorithmic}
\end{algorithm}

\section*{Dry Run Example}
Consider the one‑dimensional quadratic
\[
f(w)=(w-3)^2 .
\]
We have $\nabla f(w)=2(w-3)$, $\nabla^2 f(w)=2$.
Take $w_0=0$, stepsize $\alpha=0.1$, regulariser $\lambda=0$,
and assume *exact* gradients/Hessians (i.e. $g_t=\nabla f(w_t)$,
$H_t=\nabla^2 f(w_t)$).

\begin{center}
\begin{tabular}{c|c|c|c}
$t$ & $g_t$ & $H_t^{-1}g_t$ & $w_{t+1}=w_t-\alpha H_t^{-1}g_t$ \\ \hline
0 & $2(0-3)=-6$ & $(-6)/2=-3$ & $0-0.1(-3)=0.3$ \\
1 & $2(0.3-3)=-5.4$ & $-5.4/2=-2.7$ & $0.3-0.1(-2.7)=0.57$ \\
2 & $2(0.57-3)=-4.86$ & $-4.86/2=-2.43$ & $0.57-0.1(-2.43)=0.813$ \\
\end{tabular}
\end{center}
Corresponding objective values:
$f(0)=9$, $f(0.3)=7.29$, $f(0.57)=5.96$, $f(0.813)=4.71$.
The iterates move towards the minimiser $w^\star=3$.

\newpage
\section*{Assumptions}
\begin{assumption}[Smoothness]
$f$ has $L$‑Lipschitz gradient:
\[
\|\nabla f(x)-\nabla f(y)\|\le L\|x-y\|,\qquad\forall x,y\in\RR^{d}.
\tag{A1}
\]
\end{assumption}
\begin{assumption}[Bounded Hessian]
The true Hessian is uniformly bounded in operator norm:
\[
\|\nabla^2 f(x)\|\le M,\qquad\forall x\in\RR^{d}.
\tag{A2}
\]
\end{assumption}
\begin{assumption}[Stochastic Gradient]
\begin{enumerate}
\item Unbiasedness: $\EE[g_t\mid x_t]=\nabla f(x_t)$.
\item Bounded variance: $\EE\bigl[\|g_t-\nabla f(x_t)\|^{2}\mid x_t\bigr]\le\sigma^{2}$.
\end{enumerate}
\tag{A3}
\end{assumption}
\begin{assumption}[Stochastic Hessian]
\begin{enumerate}
\item Unbiasedness: $\EE[H_t\mid x_t]=\nabla^{2}f(x_t)$.
\item Bounded second‑moment (operator norm):
\[
\EE\bigl[\|H_t-\nabla^{2}f(x_t)\|^{2}\mid x_t\bigr]\le\tau^{2}.
\]
\end{enumerate}
\tag{A4}
\end{assumption}
\begin{assumption}[Stepsize and Regulariser]
The stepsize $\alpha$ and regulariser $\lambda$ satisfy
\[
0<\alpha\le\frac{1}{L+\lambda},\qquad \lambda>0.
\tag{A5}
\]
\end{assumption}

\section*{Main Convergence Theorem}
\begin{theorem}[Expected Gradient Norm Decay]
\label{thm:main}
Under Assumptions (A1)–(A5), let $\{x_t\}_{t=0}^{T-1}$ be generated by
\eqref{eq:update}.  Define the averaged iterate
$\bar x_T=\frac1T\sum_{t=0}^{T-1}x_t$.  Then
\[
\frac{1}{T}\sum_{t=0}^{T-1}\EE\bigl[\|\nabla f(x_t)\|^{2}\bigr]
\;\le\;
\frac{2\bigl(f(x_0)-f^{*}\bigr)}{\alpha T}
\;+\;
\underbrace{C_{1}\sigma^{2}+C_{2}\alpha\tau^{2}}_{\displaystyle\text{noise term}},
\tag{T1}
\]
where $f^{*}=\inf_{x}f(x)$ and the constants are
\[
C_{1}= \frac{2L}{\lambda},\qquad 
C_{2}= \frac{2M^{2}}{\lambda^{2}}.
\]
In particular, if $\sigma^{2}=\tau^{2}=0$ (exact gradients and
Hessians) the bound reduces to the deterministic rate
$\mathcal{O}\!\bigl(1/(\alpha T)\bigr)$.
\end{theorem}

\section*{Theoretical Properties}
\begin{itemize}
\item \textbf{Stability.}  The regulariser $\lambda$ guarantees that
  $H_t+\lambda I\succeq\lambda I$, which prevents exploding
  Newton steps even when the stochastic Hessian is indefinite.
\item \textbf{Hyper‑parameter sensitivity.}
  The bound \eqref{T1} shows a linear dependence on $\alpha$ in the
  stochastic Hessian noise term $C_{2}\alpha\tau^{2}$, suggesting that
  a too large stepsize amplifies Hessian variance, while a too small
  stepsize slows deterministic progress (the $1/(\alpha T)$ term).
\item \textbf{Comparison to first‑order methods.}
  For SGD the standard bound is $\mathcal{O}\bigl(1/\sqrt{T}\bigr)$ on the
  average gradient norm.  When the Hessian noise $\tau^{2}$ is moderate,
  SSN attains the faster $\mathcal{O}(1/T)$ deterministic rate up to an
  additive variance floor.
\end{itemize}

\section*{Discussion}
The theorem demonstrates that, despite the non‑convexity of $f$, the
subsampled Newton step yields a controlled decrease of the expected
gradient norm.  The result hinges on the unbiasedness of both stochastic
gradient and Hessian, and on the regularisation $\lambda$ that keeps the
inverse well‑behaved.  In practice, $\lambda$ may be chosen adaptively
(e.g. Levenberg–Marquardt style) to balance curvature information and
numerical stability.

\newpage
\section*{Preliminaries (Definitions and Lemmas)}
\begin{lemma}[Descent Lemma for $L$‑smooth functions]
\label{lem:descent}
For any $x,y\in\RR^{d}$,
\[
f(y)\le f(x)+\ip{\nabla f(x)}{y-x}
+\frac{L}{2}\|y-x\|^{2}.
\tag{L1}
\]
\end{lemma}
\begin{proof}
Standard; see e.g. Nesterov (2004). \qedhere
\end{proof}

\begin{lemma}[Bound on the inverse of regularised stochastic Hessian]
\label{lem:inverse}
Under (A2) and (A5) we have
\[
\lambda I\preceq H_t+\lambda I\preceq (M+\lambda)I,
\qquad
\|(H_t+\lambda I)^{-1}\|\le\frac{1}{\lambda},
\tag{L2}
\]
and
\[
\|(H_t+\lambda I)^{-1}\|_{F}^{2}\le\frac{d}{\lambda^{2}}.
\]
\end{lemma}
\begin{proof}
By (A2), $\|H_t\|\le M$ a.s.; adding $\lambda I$ shifts the spectrum
by $\lambda$, giving the eigenvalue bounds.  Inverting reverses the
inequalities, yielding the norm bound. \qedhere
\end{proof}

\section*{Main Proof (Step‑by‑Step Derivation)}
We prove Theorem~\ref{thm:main}.  Throughout, expectations are taken
conditionally on the past, i.e. $\EE_t[\cdot]=\EE[\cdot\mid x_t]$.

\bigskip
\noindent\textbf{Step 1 (Apply descent lemma).}  
Using Lemma~\ref{lem:descent} with $y=x_{t+1}$ and $x=x_t$,
\begin{align}
f(x_{t+1})
&\le f(x_t)+\ip{\nabla f(x_t)}{x_{t+1}-x_t}
+\frac{L}{2}\|x_{t+1}-x_t\|^{2}. \tag{S1}
\end{align}
Insert the update \eqref{eq:update}:
\[
x_{t+1}-x_t=-\alpha (H_t+\lambda I)^{-1}g_t.
\]

\noindent\textbf{Step 2 (Expand inner product).}
\begin{align}
\ip{\nabla f(x_t)}{x_{t+1}-x_t}
&= -\alpha\,\ip{\nabla f(x_t)}{(H_t+\lambda I)^{-1}g_t} \nonumber\\
&= -\alpha\,\ip{\nabla f(x_t)}{(H_t+\lambda I)^{-1}\nabla f(x_t)}
   -\alpha\,\ip{\nabla f(x_t)}{(H_t+\lambda I)^{-1}(g_t-\nabla f(x_t))}.
\tag{S2}
\end{align}

\noindent\textbf{Step 3 (Bound the quadratic term).}
Since $(H_t+\lambda I)^{-1}\succeq\frac{1}{M+\lambda}I$ and
$\preceq\frac{1}{\lambda}I$ (Lemma~\ref{lem:inverse}),
\begin{align}
\ip{\nabla f(x_t)}{(H_t+\lambda I)^{-1}\nabla f(x_t)}
&\ge \frac{1}{M+\lambda}\|\nabla f(x_t)\|^{2}, \tag{S3}\\
\ip{\nabla f(x_t)}{(H_t+\lambda I)^{-1}\nabla f(x_t)}
&\le \frac{1}{\lambda}\|\nabla f(x_t)\|^{2}. \tag{S4}
\end{align}
For the purpose of a lower bound on the descent we use the
\emph{upper} inequality \eqref{S4} inside the negative sign:
\[
-\alpha\,\ip{\nabla f(x_t)}{(H_t+\lambda I)^{-1}\nabla f(x_t)}
\le -\frac{\alpha}{\lambda}\|\nabla f(x_t)\|^{2}.
\tag{S5}
\]

\noindent\textbf{Step 4 (Control the stochastic cross term).}
Conditional expectation of the second term in \eqref{S2} is zero:
\[
\EE_t\bigl[\ip{\nabla f(x_t)}{(H_t+\lambda I)^{-1}(g_t-\nabla f(x_t))}\bigr]
= \ip{\nabla f(x_t)}{(H_t+\lambda I)^{-1}\EE_t[g_t-\nabla f(x_t)]}=0,
\tag{S6}
\]
by unbiasedness (A3).  Its variance will be bounded later.

\noindent\textbf{Step 5 (Bound the step norm).}
\[
\|x_{t+1}-x_t\|^{2}
= \alpha^{2}\|(H_t+\lambda I)^{-1}g_t\|^{2}
\le \alpha^{2}\|(H_t+\lambda I)^{-1}\|^{2}\|g_t\|^{2}
\le \frac{\alpha^{2}}{\lambda^{2}}\|g_t\|^{2},
\tag{S7}
\]
using Lemma~\ref{lem:inverse}.

\noindent\textbf{Step 6 (Take conditional expectation of \eqref{S1}).}
Combine \eqref{S1}, \eqref{S5}, \eqref{S7} and the fact that the cross
term vanishes in expectation \eqref{S6}:
\begin{align}
\EE_t\bigl[f(x_{t+1})\bigr]
&\le f(x_t) -\frac{\alpha}{\lambda}\|\nabla f(x_t)\|^{2}
    +\frac{L}{2}\frac{\alpha^{2}}{\lambda^{2}}\EE_t\!\bigl[\|g_t\|^{2}\bigr].
\tag{S8}
\end{align}
Decompose $\|g_t\|^{2}
= \|\nabla f(x_t)\|^{2}+\|g_t-\nabla f(x_t)\|^{2}
+2\ip{\nabla f(x_t)}{g_t-\nabla f(x_t)}$.
Taking expectation and using unbiasedness, the mixed term disappears,
yielding
\[
\EE_t\!\bigl[\|g_t\|^{2}\bigr]
= \|\nabla f(x_t)\|^{2}+\EE_t\!\bigl[\|g_t-\nabla f(x_t)\|^{2}\bigr]
\le \|\nabla f(x_t)\|^{2}+\sigma^{2},
\tag{S9}
\]
by variance bound (A3).

\noindent\textbf{Step 7 (Insert \eqref{S9} into \eqref{S8}).}
\begin{align}
\EE_t\bigl[f(x_{t+1})\bigr]
&\le f(x_t) -\frac{\alpha}{\lambda}\|\nabla f(x_t)\|^{2}
    +\frac{L\alpha^{2}}{2\lambda^{2}}\Bigl(\|\nabla f(x_t)\|^{2}+\sigma^{2}\Bigr) \nonumber\\
&= f(x_t) -\underbrace{\Bigl(\frac{\alpha}{\lambda}
    -\frac{L\alpha^{2}}{2\lambda^{2}}\Bigr)}_{\displaystyle\gamma}
    \|\nabla f(x_t)\|^{2}
    +\frac{L\alpha^{2}\sigma^{2}}{2\lambda^{2}} .
\tag{S10}
\end{align}
Because of the stepsize restriction (A5), $\alpha\le 1/(L+\lambda)$,
we have $\gamma\ge \dfrac{\alpha}{2\lambda}>0$.
Thus,
\[
\EE_t\bigl[f(x_{t+1})\bigr]
\le f(x_t)-\frac{\alpha}{2\lambda}\|\nabla f(x_t)\|^{2}
    +\frac{L\alpha^{2}\sigma^{2}}{2\lambda^{2}} .
\tag{S11}
\]

\noindent\textbf{Step 8 (Account for Hessian stochasticity).}
The bound \eqref{S11} used only Lemma~\ref{lem:inverse}, which relies on
a deterministic spectral bound $M$.  Since $H_t$ itself is random,
we must add the error caused by the deviation
$\Delta_t:=H_t-\nabla^{2}f(x_t)$.
Rewrite the update as
\[
x_{t+1}=x_t-\alpha\bigl(\nabla^{2}f(x_t)+\lambda I+\Delta_t\bigr)^{-1}g_t .
\]
Using the matrix identity
\[
(A+\Delta)^{-1}=A^{-1}-A^{-1}\Delta A^{-1}+A^{-1}\Delta A^{-1}\Delta (A+\Delta)^{-1},
\tag{S12}
\]
with $A:=\nabla^{2}f(x_t)+\lambda I$, and taking norms,
\[
\bigl\|(A+\Delta)^{-1}-A^{-1}\bigr\|
\le \|A^{-1}\|^{2}\|\Delta\|
    +\|A^{-1}\|^{3}\|\Delta\|^{2}.
\tag{S13}
\]
Because $\|A^{-1}\|\le 1/\lambda$ and $\|\Delta\|$ has bounded second
moment $\tau^{2}$ (Assumption A4), we obtain
\[
\EE_t\!\bigl[\|(H_t+\lambda I)^{-1}-(\nabla^{2}f(x_t)+\lambda I)^{-1}\|^{2}\bigr]
\le \frac{2\tau^{2}}{\lambda^{4}} .
\tag{S14}
\]

When expanding the descent inequality we encounter the term
\[
\alpha\,\ip{\nabla f(x_t)}{[(H_t+\lambda I)^{-1}-(\nabla^{2}f(x_t)+\lambda I)^{-1}]\,g_t } .
\]
Using Cauchy–Schwarz, boundedness of $\|g_t\|$ (from (S9)) and
\eqref{S14},
\[
\EE_t\bigl|\text{Hessian error term}\bigr|
\le \alpha\,
\frac{1}{\lambda}\|\nabla f(x_t)\|
\cdot\frac{\sqrt{2}\,\tau}{\lambda^{2}}\bigl(\|\nabla f(x_t)\|+\sigma\bigr)
\le \frac{\alpha\sqrt{2}\,M}{\lambda^{3}}\|\nabla f(x_t)\|^{2}
    +\frac{\alpha\sqrt{2}\,\tau\sigma}{\lambda^{3}} .
\tag{S15}
\]
Absorbing the first component into the coefficient of
$\|\nabla f(x_t)\|^{2}$ (it merely reduces $\gamma$) and the second
component into a constant yields an additional term
$\frac{C_{2}\alpha\tau^{2}}{2}$ with $C_{2}=2M^{2}/\lambda^{2}$.

\noindent\textbf{Step 9 (Full one‑step expectation inequality).}
Collecting deterministic descent \eqref{S11} and the Hessian‑noise
contribution \eqref{S15}, we obtain
\[
\EE_t\bigl[f(x_{t+1})\bigr]
\le f(x_t)-\frac{\alpha}{2\lambda}\|\nabla f(x_t)\|^{2}
    +\frac{L\alpha^{2}\sigma^{2}}{2\lambda^{2}}
    +\frac{C_{2}\alpha\tau^{2}}{2}.
\tag{S16}
\]

\noindent\textbf{Step 10 (Unroll the recursion).}
Take total expectation and sum \eqref{S16} from $t=0$ to $T-1$:
\begin{align}
\EE\bigl[f(x_T)\bigr]
&\le f(x_0)-\frac{\alpha}{2\lambda}\sum_{t=0}^{T-1}
      \EE\!\bigl[\|\nabla f(x_t)\|^{2}\bigr]
   +T\Bigl(\frac{L\alpha^{2}\sigma^{2}}{2\lambda^{2}}
          +\frac{C_{2}\alpha\tau^{2}}{2}\Bigr).
\tag{S17}
\end{align}
Since $f(x_T)\ge f^{*}$, rearrange:
\[
\frac{1}{T}\sum_{t=0}^{T-1}\EE\!\bigl[\|\nabla f(x_t)\|^{2}\bigr]
\le
\frac{2\lambda\bigl(f(x_0)-f^{*}\bigr)}{\alpha T}
   +\frac{L\alpha\sigma^{2}}{\lambda}
   +C_{2}\tau^{2}.
\tag{S18}
\]
Define $C_{1}=2L/\lambda$ and recall $C_{2}=2M^{2}/\lambda^{2}$, then
\eqref{S18} is exactly the bound \eqref{T1} of Theorem~\ref{thm:main}.
\qed

\section*{Rate Derivation (With Constants)}
From \eqref{T1} we see three contributions:
\begin{enumerate}
\item \textbf{Deterministic decay term}
  $\displaystyle\frac{2(f(x_0)-f^{*})}{\alpha T}$,
  which vanishes as $T\to\infty$ at rate $\mathcal{O}(1/T)$.
\item \textbf{Gradient‑noise floor}
  $C_{1}\sigma^{2}= \frac{2L}{\lambda}\sigma^{2}$,
  independent of $T$.
\item \textbf{Hessian‑noise floor}
  $C_{2}\alpha\tau^{2}= \frac{2M^{2}}{\lambda^{2}}\alpha\tau^{2}$,
  linear in $\alpha$.
\end{enumerate}
Choosing $\alpha$ to balance the two variance floors gives
\[
\alpha^{\star}= \min\Bigl\{\frac{1}{L+\lambda},
      \sqrt{\frac{L\lambda}{M^{2}}}\,\frac{\sigma}{\tau}\Bigr\},
\]
which minimises the right‑hand side of \eqref{T1}.

\section*{Special Cases}
\begin{itemize}
\item \textbf{Exact gradients, noisy Hessian ($\sigma=0$, $\tau>0$).}
  The bound reduces to
  $\displaystyle\frac{1}{T}\sum_{t}\EE\|\nabla f(x_t)\|^{2}
   \le \frac{2\lambda(f(x_0)-f^{*})}{\alpha T}
   +\frac{2M^{2}}{\lambda^{2}}\alpha\tau^{2}$.
  Taking $\alpha\asymp 1/\sqrt{T}$ yields a $\mathcal{O}(1/\sqrt{T})$ rate.
\item \textbf{Exact Hessian, noisy gradient ($\tau=0$, $\sigma>0$).}
  We recover the classical SGD‑type bound
  $\displaystyle\frac{2\lambda(f(x_0)-f^{*})}{\alpha T}
   +\frac{2L}{\lambda}\sigma^{2}$.
\item \textbf{Both exact ($\sigma=\tau=0$).}
  The algorithm achieves the deterministic Newton rate
  $\mathcal{O}(1/(\alpha T))$, i.e. linear convergence for strongly
  convex regions (not covered here) and $\mathcal{O}(1/T)$ for general
  smooth non‑convex functions.
\end{itemize}

\end{document}